\begin{center}
    \section*{Abstract}
\end{center}
Fire incidents in densely populated urban areas in Bangladesh result in substantial loss of life and property. Existing alarm systems often fail to provide rapid detection, robust communication, and coordinated emergency response. This thesis introduces an integrated fire-response framework that incorporates autonomous fire and smoke detection, network-independent alert relays, acoustic-wave fire suppression, and coordinated emergency dispatch. The detection layer employs a YOLOv8n model with a heuristic fallback, while the communication layer integrates SIP-over-UDP voice calls with an OpenWrt-based peer-to-peer mesh network to ensure alert delivery during network disruptions. A low-frequency acoustic suppressor is implemented to extinguish small flames, and a Next.js dispatch portal facilitates station assignment, crew and vehicle tracking, escalation, and citizen SOS reporting. The system was evaluated using 73,600 fire and smoke images from the BD-Threats-v1 corpus. The YOLOv8n model achieved a mean average precision at 0.5 (mAP@0.5) of 0.91 and processed 28 frames per second on a Raspberry Pi 4. Detection-to-alert latency was measured at 410 milliseconds with an active uplink and 1.8 seconds via mesh relay. The acoustic suppressor extinguished a 15-centimeter alcohol pan fire in 4.6 seconds at 45 Hz. Collectively, the proposed framework establishes a practical foundation for resilient and proactive urban fire-response systems.
\vspace{5mm}

\textit{Keywords: Fire Detection, YOLOv8n, Mesh Networking, Acoustic Fire Suppression, Emergency Dispatch}